\section{Introduction}
\label{sec:intro}

A vector database stores items such as text passages as vectors.
A vector is a list of numbers that an AI language model makes from a passage.
Think of each vector as a point on a map, where passages with similar meaning sit close together.
A search question, or query, becomes a vector the same way.
The database returns the stored passages closest to it, which we call the query's best matches.

A distributed vector database splits its data into parts called shards, each kept on its own server.
A central server, the coordinator, sends each query to some of the shards and merges their answers.
It works like a library with many branch buildings and one front desk.
In the system we study, the coordinator sends each query to only two of the \RNumNumShards{} shards.
To choose them, it keeps one summary per shard: a short list of vectors that mark where its passages lie.
A summary is like a catalog card that lists the main topics a branch holds.
It picks the two whose summaries come closest to the query.

New passages keep arriving, but the coordinator keeps each shard's old summary until that shard sends a fresh one.
An old summary is like an old catalog card that does not list the newest books.
So a query may skip the shard that holds its best new match, and miss that match (Fig.~\ref{fig:problem}).
This matters for questions about newly added documents, because their answers are in the newest passages.

\begin{figure}[!t]
\centering
% Figure 1 drawing (TikZ). Every label is a word used in the text.
\begin{tikzpicture}[yscale=0.79, font=\scriptsize, >={Stealth[length=4pt]},
  shard/.style={draw, rounded corners=2pt, minimum width=1.15cm, minimum height=0.8cm, fill=white, align=center},
  summ/.style={draw, minimum width=0.42cm, minimum height=0.42cm, inner sep=0pt, fill=white},
  tag/.style={circle, draw, very thick, fill=white, inner sep=0.6pt, font=\scriptsize\bfseries}]
% Coordinator with one summary per shard and a store for pushed passages.
\draw[rounded corners=3pt, fill=black!5] (-4.2,2.15) rectangle (4.2,4.0);
\node[anchor=north west, font=\footnotesize\bfseries] at (-4.15,3.97) {Coordinator};
\node[anchor=north west] at (-4.15,3.55) {one summary per shard (\RNumCentersPerSummary{} vectors each)};
\node[summ] (u1) at (-3.8,2.75) {1};
\node[summ] (u2) at (-3.2,2.75) {2};
\node[summ, draw=red!75!black, dashed, very thick] (u3) at (-2.6,2.75) {3};
\node at (-2.0,2.75) {$\cdots$};
\node[summ] (u32) at (-1.4,2.75) {32};
\node[anchor=west, text=red!75!black, align=left] at (-1.05,2.75) {summary 3 is\\out of date};
\node[draw, fill=white, align=center] (pushed) at (2.9,2.85) {pushed passages,\\[-1pt]checked for\\[-1pt]every query};
% Query: it enters the coordinator, which compares it with the summaries and also checks the pushed passages.
\node[anchor=south] (q) at (1.0,4.35) {query};
\draw[->, thick] (q.south) to (1.0,3.38) to (0.32,3.38);
\draw[->, thick, densely dotted] (1.0,3.38) to node[above, pos=0.6, inner sep=1.5pt] {also checks} (pushed.north);
% Shards.
\node[shard] (s1) at (-3.45,0) {Shard 1};
\node[shard] (s2) at (-2.05,0) {Shard 2};
\node[shard, minimum width=2.1cm] (s3) at (0.15,0) {Shard 3\\[1pt]{\color{red!75!black}\hspace{9pt}new match}};
\node[star, star points=5, star point ratio=2.3, fill=red!75!black, inner sep=1.15pt] at (-0.47,-0.15) {};
\node at (1.9,0) {$\cdots$};
\node[shard] (s32) at (3.45,0) {Shard 32};
% The two searched shards.
\draw[->, thick] (-3.45,2.15) to (-3.45,0.4);
\draw[->, thick] (-2.05,2.15) to (-2.05,0.4);
\node at (-2.75,1.25) {searched};
% New passages arriving at shard 3.
\foreach \x in {-0.2,0.05,0.3,0.55} { \fill[red!75!black] (\x,-0.8) circle (1.6pt); }
\draw[->, red!75!black, thick] (0.15,-0.7) to (0.15,-0.4);
\node[anchor=north, text=red!75!black] at (0.15,-0.9) {new passages};
% The three options: arrows with number tags.
\draw[->, very thick, blue!70!black] (-0.65,0.4) to node[pos=0.5, tag, draw=blue!70!black, text=blue!70!black] {1} (u3.south);
\draw[->, very thick, green!45!black, dashed] (0.4,2.15) to node[pos=0.5, tag, draw=green!45!black, text=green!45!black] {2} (0.4,0.4);
\draw[->, very thick, orange!85!black] (1.05,0.4) to node[pos=0.5, tag, draw=orange!85!black, text=orange!85!black] {3} (pushed.south);
% Key for the three options.
\node[anchor=north west, align=left] at (-4.15,-1.35) {%
{\color{blue!70!black}\textbf{(1) refresh:}} a shard sends a fresh summary\\
{\color{green!45!black}\textbf{(2) extra search:}} the coordinator also asks the shard whose summary ranks third\\
{\color{orange!85!black}\textbf{(3) push:}} shards send copies of new passages to the coordinator};
\end{tikzpicture}

\caption{The problem and the three options. Shard~3 has new passages, but its summary at the coordinator is out of date, so the query can skip shard~3 and miss a new match. Shard~3's summary ranks third closest to the query, so extra search would also ask it.}
\label{fig:problem}
\end{figure}

We study this problem with a simulator, a program that imitates the system one time step at a time.
Each run lasts \RNumNumEpochs{} time steps, and about \RNumSteadyEpochInsertVecs{} new passages arrive per step on average.
In each time step, the coordinator also answers a batch of queries, \RNumPlainQueriesDefault{} by default.
The system may send upkeep messages, paid from a fixed budget of network bytes per time step.
The budget works like a monthly phone data plan without rollover: unused bytes are lost.
We count it in summaries: our default budget of four pays for four summaries per step.
We also test budgets of two and eight.
Even the largest is too small to refresh every shard each step.
So the question is how to spend the budget, and we compare three options (Fig.~\ref{fig:problem}).
The first option is to refresh: some shards send fresh summaries, chosen by a refresh rule.
One such rule, most-changed, picks the shards whose data changed most.
The second is extra search: when the coordinator is unsure which shards to ask, it also asks a third one.
The third is push: shards send copies of their newest passages to the coordinator, which checks them for every query.

We measure accuracy with recall, the share of a query's true best matches that the system returns.
For example, if the system returns seven of ten true best matches, its recall is seven in ten.
Recent recall counts only the best matches that arrived in the last two time steps.
It shows how well the system finds new data, while overall recall counts all best matches, old or new.
When we say that a method wins or is better, we mean that it has higher recent recall.
A difference is clear when our statistical test shows that chance alone would rarely produce it.
A clear lead is fragile when a further check casts doubt on it (Section~\ref{sec:stats}).
Our simulator counts the bytes of every message the budget pays for.
It checks every answer against the exact best matches.

We vary where new passages land: with the most similar shard, all on one shard, or spread over eight shards.
New passages either mix topics the shards already hold or bring new topics, one at a time.
We wrote down our plan and tests before the main analysis; Section~\ref{sec:setup} lists its changes.

When each new passage joins the most similar shard, old summaries often still send queries there.
Refreshing then gains little, and at the default budget and query rate extra search and push beat every refresh rule.
When one shard takes all of a time step's new passages, refresh wins by far on new topics.
On mixed topics, fresh summaries barely help with new matches and slightly lower accuracy overall, so push wins.
When eight shards share the new passages, push wins on mixed topics.
On new topics, refresh has only a fragile lead at a budget of two.
Neither refresh nor push clearly wins at four, and push wins at eight.

Refresh raises a second question: which shards should refresh?
Our sample-query probe, one of five refresh rules, uses practice queries to estimate which fresh summaries would help most.
With eight shards on new topics, it clearly lost to most-changed at budgets of four and eight.
The probe's own messages cost it one summary, but most-changed's small score messages were free.
So we ran a follow-up in which most-changed pays for them, chosen after seeing the main results and so exploratory.
It removed the fragile lead over push.
It also showed that at budgets of four and eight, the probe lost mostly by picking worse shards.
On average, counting each shard's new passages predicted the value of a fresh summary better than the probe did.
We did not try that count as a refresh rule.
